\subsection{Overview}

Despite decades of field experimentation and steady advances in numerical
simulation, the community that studies water movement in subsurface (tile)
drainage systems consistently identifies a set of persistent, mutually
reinforcing knowledge gaps. Across studies spanning humid temperate cropland,
arid and semi-arid irrigated plains, saline coastal soils, and tropical
oases, authors converge on a common message: the processes that most strongly
govern drain hydrographs---preferential and macropore flow, spatially and
temporally variable soil hydraulic properties, and the coupling between soil
pore-water flow, pipe flow, and surface water---remain incompletely
characterised and imperfectly represented in the models used to guide design.
This section frames the research frontier around five interlocking lines of
inquiry that recur throughout the literature: (i) the characterisation and
mechanistic modelling of preferential and macropore flow; (ii) the need for
coupled surface--subsurface and pipe-flow hydraulic models, including
pressurised-flow and open-ditch processes; (iii) the quantification of
spatial and temporal heterogeneity in soil hydraulic properties and the data
required to constrain it; (iv) improved calibration, validation, and
long-term monitoring protocols; and (v) the revision of design criteria and
management strategies, including controlled drainage, novel drainage
structures, and adaptation to climate change. A cross-cutting theme is that
progress in each area is limited less by conceptual understanding than by the
availability of field data at the right spatial and temporal resolution to
test, calibrate, and generalise the models. Headline priorities include
mechanistic characterisation of macropore--drain connectivity, coupling of
in-pipe surface flow with soil pore-water flow, denser and depth-resolved
hydraulic-conductivity sampling, longer continuous monitoring records, and the
explicit incorporation of climate-change scenarios into design.

\subsection{Preferential and macropore flow: characterisation and modelling}

The representation of preferential flow is perhaps the most frequently cited
frontier. \textcite{Gaillot2023} argue that the relative contribution of
preferential flow to drains needs to be better characterised as a function of
soil type, soil water content, and agricultural practice, and that models must
then be tested for extrapolation to other fields while accounting for soil
variability and drainage-network characteristics. \textcite{Fox2004}
similarly call for the degree of direct connectivity between macropores and
subsurface drains to be characterised using more mechanistic approaches, for
the ``express fraction'' to be evaluated under varying saturated--unsaturated
flow conditions above the tile drain, and for the modified model to be tested
under different climate, soil, drain-configuration, and chemical-application
conditions. In the U.S. Midwest, \textcite{Askar2020} recommend further
testing of the simple empirical approaches DRAINMOD uses to represent
macropore flow, emphasising that rigorous testing requires water-table-depth
measurements alongside drainage and surface-runoff data, direct field
measurement of macropore-flow inputs, and detailed data on crack geometry,
temporal crack dynamics with soil moisture, and the area contributing to
surface runoff.

Improvements to macropore process descriptions are proposed across several
model frameworks. \textcite{Li2023d} call for further study of preferential
flow in macropore channels and appropriate modification of the RZWQM2
macropore component---better partitioning of soil water between matrix and
macropore flow, improved macropore channel flow to shorten the drainage-peak
recession time, and calibration of macropore parameters against sub-hourly
drainage data. At watershed scale, \textcite{Poon2021} suggest adjusting the
percolation algorithm to account for dynamic macropore connectivity or
effective macroporosity and to include a soil-moisture factor describing
matrix--macropore water transfer during the growing season, noting that
experimental measurements of soil macroporosity and site-specific measurements
of major soil types and management practices are needed to select appropriate
connectivity values. \textcite{Villholth1998b} propose including small-scale
diffusion across the micropore space or adding more domains to the model
description, while cautioning that such remedies demand more elaborate models
and more parameters, and stress the importance of site-specific measurements
of matrix and macropore hydraulic conductivity. \textcite{Villholth1998}
further note that the dynamics of macropore structure at the soil surface and
their influence on chemical transport deserve additional study, and that
whether vertical macropore transport occurs in deeper saturated layers---
especially in areas dominated by regional groundwater recharge---remains to be
investigated.

Several groups highlight structural and dynamic controls on macroporosity that
models neglect. \textcite{Holbak2021} recommend including hysteresis effects,
investigating the interrelationship between biopore-flow initiation and
hysteresis under laboratory and field conditions, and validating biopore
models across the soil profile against datasets combining biopore
measurements, soil-water dynamics, and solute leaching in drained fields.
\textcite{Zhao2001} raise the possibility that additional years of manure
application could significantly improve soil structure and macropore
development, and \textcite{Wang2022} identify the reestablishment of
macropores and nutrient transfer through cracks as targets for improving EPIC
simulation precision. \textcite{Varvaris2018} call for scrutiny of the
calibrated saturated hydraulic conductivity of macropores---and the resulting
indirect estimate of potential preferential-pathway flow---and propose a
three-dimensional dual-permeability model to capture spatial geological
layering, lateral macropore-flow velocities, and non-systematically
tile-drained fields. Tracer- and geophysics-based identification of
preferential pathways is repeatedly recommended: \textcite{Tähtikarhu2025}
suggest ground-penetrating radar and tracers for flowpath identification and
the inclusion of preferential flowpaths, more comprehensive open-ditch
processes, and soil-frost processes to improve model representation, while
\textcite{Salo2017} argue that direct measurements of flow routes are more
useful than aggregated discharge for calibrating spatially variable parameters
and that tracer-tracked pathways require solute-transport simulation.

\subsection{Coupled surface--subsurface, pipe-flow, and open-ditch hydraulics}

A recurrent structural gap is the absence of models that couple soil
pore-water flow with surface flow inside drains and ditches.
\textcite{Xin2016b} call for future model development to couple surface water
flow in drains with pore-water flow in soils to represent the temporally
varying hydraulic-head condition around the drain, to extend analytical
solutions to include rainfall and evapotranspiration via source and sink terms
for combined drainage and infiltration, and to develop a more comprehensive
treatment of heterogeneous soils, complex pore-water flow, and uncertainty in
spatially variable capillary-fringe thickness and capillary rise.
\textcite{Nousiainen2015} propose enabling simulation of water flow in drain
pipes within FLUSH to capture pressurised flow and clogging effects, adding a
description of open-ditch networks to extend the model to larger areas, and
assigning distinct potential-evapotranspiration series to areas under
different crops. Pressurised-flow dynamics are a specific frontier:
\textcite{Henine2010} argue that coupling subsurface pipes with a feedback
linkage to the water table could introduce a validated temporary-retention
function under conditions exceeding design capacity and broaden the scope of
drainage-hydrology investigations, and \textcite{Henine2014} add that a
complete study of drainage overpressure at the catchment outlet should account
for surface runoff and its interaction with the water table during pipe
pressurisation. \textcite{Saadat2018} similarly suggest separately analysing
periods when pipes are pressurised using pressure-transducer data.

Entrance resistance and near-drain flow representation remain unresolved for
older and heterogeneous systems. Entrance
resistance must be considered when modelling old drainage systems---for
example by assigning a higher pressure head to water in gravel at the drain
boundary---and a smaller element size is needed near the drain for
detailed pressure-head study. \textcite{Messing2006} raise the question of
whether heavy-rain water transport should be characterised by a rapid onset of
flow through the permeable backfill trench followed by slower, continuous
transport from the inter-trench zone, and how this partitioning influences
drain-water quality. The choice of model dimensionality is itself a research
question: \textcite{Alakukku2010} advise a two- or three-dimensional approach
when predicting the proportional fractions of drain flow and surface runoff or
simulating groundwater level in sloping fields, and \textcite{Carlier2007}
call for further development before applying their method to larger domains
such as networks of drain plots, simpler modelling of drainage collectors, a
better representation of soils with extreme hydrodynamic parameters, and a more
precise averaging of hydraulic conductivities for unsaturated flow.
\textcite{FIPPS1989} note that the regions at the top and bottom of a slope,
and the entire flow region when the water table falls below the drains, require
additional analysis. Broader integrating frameworks are also sought:
\textcite{Manguerra1998} argue that SWAT should be extended to include and
directly handle the physical processes governing water movement to subsurface
drains, expecting better performance from explicit representation, and
\textcite{SAVABI1993} report that subsurface lateral water movement and
rainfall interception---absent from the WEPP hydrology model---are being
incorporated, alongside a need for better hydraulic-roughness estimation for
runoff routing.

\subsection{Spatial and temporal heterogeneity of soil hydraulic properties}

Nearly every modelling study attributes residual error to unresolved
heterogeneity in soil hydraulic properties, particularly hydraulic
conductivity. \textcite{DeSchepper2017} suggest that a more detailed
representation of macropores, fractures, sand lenses, and other
heterogeneities could improve reproduction of observed hydraulic-head
variability, and that including specific storage and residual saturation in
calibration would improve performance. \textcite{Wesström2003} note that other
approaches are needed to account for variation of properties such as hydraulic
conductivity with depth, and \textcite{Karandish2018} propose increasing the
number of soil layers or gradually varying hydraulic parameters, adopting a
dual-porosity model for paddy soils, and accounting for temporal variation in
hydraulic parameters during simulation. Several authors call explicitly for
denser and depth-resolved sampling: \textcite{PalaciosVélez2004} recommend at
least one hydraulic-conductivity measurement per hectare (five per plot for
plots under 5~ha) and a more detailed study of the spatial distribution of
conductivity, the location of the impermeable layer, and the appropriate
drain-spacing approach, together with post-construction comparison of designed
and observed conductivity and water-table height. \textcite{Messing2006} argue
that more points along the transect between trench lines are needed to capture
spatial variability, and \textcite{Tähtikarhu2025} call for more measurements
of vertical hydraulic conductivity, measurement of horizontal conductivities,
and a more sophisticated description of water-retention properties that
accounts for hydrophobicity and hysteresis at field scale.

Anisotropy and long-term change in soil properties form a related sub-frontier.
\textcite{Akis2015} call for studies of long-term spatial variability under
drained and undrained conditions with anisotropy, a vadose-zone hydrological
model to understand salt-movement and water-flow relations to saturated
conductivity, and understanding of spatial anisotropy in the drainage--
infiltration relationship in the upper 1.5~m. \textcite{Moustafa1998} likewise
call for study of long-term spatial variation of soil-water properties as
affected by subsurface drainage and closer-spaced sampling to resolve true
spatial structure. \textcite{Koivusalo1999} recommend relating soil hydraulic
conductivity to root-zone deficit, describing time-dependent soil-structure
changes during drying, and obtaining measured field-scale values for
macroporous clayey soil. \textcite{DeBoer2006} stress that saturated hydraulic
conductivity and the relationship between drainage volume and upward flux must
be defined with reasonable accuracy, supported by a comprehensive soil survey
as DRAINMOD input. At watershed scale, \textcite{Schilling2015} suggest
incorporating subsoil porosity variations by linking porosity to SSURGO
bulk-density data and spatially refining model results with greater horizontal
and vertical resolution of soil hydraulic properties.

\subsection{Calibration, validation, parameterisation, and data needs}

Beyond the properties themselves, authors emphasise improved calibration
strategy and the acquisition of complementary observations.
\textcite{Boico2022} recommend calibration with parameter-estimation software
to overcome poor representation of soil hydraulic properties and the
implementation of multiple conceptual models to represent plausible
heterogeneity structures. \textcite{Schneider2022} propose a more flexible
parameterisation for drain-sensitive parameters such as the conductivity of
the uppermost geologic layers, a fully distributed conductivity field with
pilot-point calibration, additional observation types such as soil moisture to
reduce overparameterisation, finer horizontal and vertical resolutions with a
more complex unsaturated-zone description and explicit drain representation,
and integration of proxy datasets such as machine-learning predictions of
drain probability or average drain flow with physically based models.
\textcite{DeSchepper2017} suggest combining catchment- and drainage-discharge
calibration, and \textcite{Sedaghatdoost2019} recommend using drainage
salinity data alongside other measured data in the objective function.
\textcite{Brighenti2023} propose testing alternative objective functions such
as the Kling--Gupta efficiency and establishing relationships between
plot-scale measured tile flow and watershed-scale calculated tile flow.

Complementary in-situ observations are repeatedly identified as limiting.
\textcite{Askar2020} stress that water-table-depth measurements must accompany
drainage and surface-runoff data. \textcite{Salo2017} argue that more output
variables such as groundwater level should be considered when assessing
drainage practices. \textcite{Varvaris2018} identify a need to determine the
unknown source continuously recharging a tile network during dry periods, to
obtain in-situ soil-water-content measurements for a more reliable calibrated
flow model, and to run conservative-tracer tests with solute-transport
simulation. \textcite{Hwang2024} call for extended field observations---soil
profiling, hydraulic conductivity, water-retention curves, infiltration rates,
water-table fluctuations, and soil-moisture change over time---and a
complementary modelling study of how drainage design influences hydrology and
crop yield. Several authors emphasise that model claims must be tested against
field data before recommendations are made: \textcite{Askar2021} state that
the potential limitations of DRAINMOD-P must be evaluated against
field-measured data prior to recommending further modifications;
\textcite{Salazar2008} caution that tests on different soil types are needed
before recommending DRAINMOD runs with ROSETTA-estimated saturated
conductivity; \textcite{Singh2020b} and \textcite{Malota2024} call for
additional testing across regions, soils, and climates; and
\textcite{Holbak2021} require further field and laboratory experiments to
validate models across the soil profile. \textcite{Sinai2006} recommend
verifying that field conditions match the assumptions underlying DRAINMOD, and
\textcite{He2002} note that models other than DRAINMOD should be used when
groundwater flow lines must be computed across a landscape.

\subsection{Long-term monitoring and revision of design criteria}

A strong current in the literature is that short-term studies mislead and that
design guidance must rest on long records. \textcite{Wang2006b} argue that
drainage design should not be based on year-by-year short-term DRAINMOD
simulations, that long-term runs give better average yield predictions to
guide spacing, and that the benefits of controlled drainage must be
demonstrated and quantified locally. \textcite{Ale2009} call for analysing
drainage water management under wider spacings (20--40~m) over an approximately
100-year climatic record and for exploring uncertainty in vertical-conductivity
calibration. \textcite{Häggblom2019} note that studying interannual
variability in supplementary drainage requires longer simulation periods with
hourly input data, and \textcite{Malota2015} state that long simulation
periods are required to assess shallow-water-table control and that verifying
design criteria is paramount before adoption. \textcite{Salo2019} recommend a
longer monitoring period to reveal long-term effects of installation method,
studying variation between sections of the same drainage method to separate
machinery effects, and automatic measurements to compute drainage response time
after rain events. Long-term monitoring and
evaluation of drainage configurations are emphasised to validate results and
refine design recommendations, and \textcite{Kladivko2023} caution against attributing
drainage responses to a management practice that occurred only shortly before.
\textcite{Chang2012} call for longer observation of water and salt movement in
deep soil layers, and \textcite{Deelstra2015} anticipate further analysis of
soil physical data and better description of the interaction between weather,
soil physical conditions, groundwater behaviour, and drain discharge as
prerequisites for correct design under climate change.

Many site studies conclude that established design criteria require
re-evaluation. \textcite{Christen2001} call for new drainage design criteria
protecting crops while minimising drain-water salinity and volume,
reassessment of criteria in light of land-use and irrigation change, and novel
saline-water disposal options. \textcite{Muhammad2021} recommend updating
criteria and decreasing drain depth to 120~cm at some sites and 100~cm at
others while holding spacing constant, and \textcite{Malota2016} call for
revisiting design criteria to account for land slope and for recalculating
drain spacing (currently 54~m in clay and 74~m in clay-loam soils).
\textcite{Bahçeci2009} argue that a lower drain-discharge rate should inform
future design in the Harran Plain, while \textcite{Kahlown2002} conclude that
subsurface drainage interventions must not go deeper than 2~m and that findings
still require field and laboratory confirmation. \textcite{Mante2022} state
that the recommended 12--15~m (up to 90~m) spacing for sandy loams must be
reviewed and paired with controlled-traffic farming at narrower spacing
($\le$12~m, 0.9~m depth), and \textcite{Nelson2017} call for determining the
cost-effectiveness of tile spacings below 6.1~m in claypan soils.
\textcite{HajAmor2019} identify a need to design new drainage systems---
including standard guidelines for depth and spacing---and to test them in situ
before extension, and \textcite{Skaggs2017} advocate standardising the
definition of the drainage coefficient and the consistent use of the Kirkham
coefficient and drainage intensity to enable comparison across studies.
\textcite{GALLICHAND1992b} suggest implementing design procedures within
CAD/GIS-based drainage-design software.

\subsection{Controlled drainage, drainage water management, and novel structures}

Controlled drainage and drainage water management (DWM) constitute a
management frontier where hydrologic understanding is still maturing.
\textcite{Youssef2021} call for well-coordinated experimental and modelling
research on how seepage affects controlled-drainage performance for reducing
flow and nitrogen losses, and for monitoring guard lines of free-drainage plots
to evaluate lateral seepage from managed to non-managed field portions;
\textcite{Pease2017} likewise cite field experiments as necessary to represent
seepage in discharge reduction and propose examining slow spring
profile-drawdown strategies. \textcite{Lahdou2019} recommend examining
potential increases in overland flow due to controlled drainage through field
measurement or detailed modelling, and \textcite{Ale2009} note the influence of
vertical-conductivity uncertainty on the DWM contribution to seepage.
\textcite{Craft2018} argue that modelling excess-moisture stress is vital to
understanding drawbacks of DWM adoption, that drainage routines within a
controlled-outlet system may need re-examination to simulate volume when the
outlet opens in spring, and that spring nitrogen loading is a priority.
A further priority is extending representative Midwest soils,
including multiple GCMs, verifying the model's ability to simulate DWM effects
on yield and drainage volume with field data, and evaluating climate-change
impacts on drainage-climate outputs. \textcite{VandeCraats2021} propose
expanding anticipatory drainage-base management to include evapotranspiration,
subirrigation, spatially variable seepage, and heterogeneity, and to address
which weather forecasts and forecast lengths to use and how forecast
uncertainty affects confidence in optimal strategies. \textcite{Salla2022}
call for investigating how collector-ditch water-level management can increase
controlled-drainage efficiency and how control schedules affect the
groundwater table seasonally and under future conditions, and
\textcite{Strock2025b} identify multiple and irregularly timed inflow sources
as a knowledge gap and advocate automation of low-grade weirs to adjust
drainage intensity responsively. \textcite{Güngör2023} caution that
controlled drainage should not be overlooked because a very rapid water-table
decrease may increase irrigation demand in dry summers, and
\textcite{Dou2023} note the need to adjust drainage depth, spacing, and period
to local characteristics.

Novel drainage structures and materials define a further set of open
questions. \textcite{Ren2023} call for improving HYDRUS-2D boundary
conditions, studying whether the negative pressure of siphon-type drainage
changes during head drawdown, verifying siphon drainage in the field, and
considering pipe perforation rate and geotextile specifications.
\textcite{Tao2025b} recommend long-term measurement to determine the working
time of their system relative to construction costs and to balance reduced
waterlogging with soil-water retention and pollution reduction, and
\textcite{Li2021} propose new structures that adjust water-holding capacity,
permeability, and aeration while eliminating capillary barriers and hysteresis,
inhibiting rapid vertical movement and bypass flow so that water converges
horizontally into the pipe. Multi-level and bi-level configurations warrant
further study: \textcite{Hornbuckle2007} call for investigating multi-level
systems under differing conditions, actively managing shallow and deep drains,
and increasing deep-drain spacing commensurate with shallow-drain drainage,
and \textcite{Verma1998} suggest bi-level drainage as controlled drainage via
the depth of the shallow drain. Envelope and perforation design remain open:
\textcite{Alavi2021} recommend testing the Hydroluis drainpipe under different
gradients, rates, soils, and inter-pipe distances; \textcite{Wu2022} call for
considering envelope-material conductivity, soil textures, and perforation
shape and pattern; \textcite{Oyarce2017} suggest installing an experimental
prototype in an actual soil profile and evaluating different screen types; and
\textcite{HajAmor2019} recommend replacing existing systems with geotextile
envelopes to prevent sand clogging. Numerical techniques can characterise fluid behaviour in larger module
arrangements and flow conditions without expensive laboratory experiments, and
\textcite{Nelson2011b} observe that a drainage/subirrigation system should have
remained in drainage mode longer than tested.

Optimal layout and burial-depth determination in saline and problem soils
recur in recent Chinese studies. \textcite{Tian2023} stress considering the
appropriate ditching depth when determining optimal pipe-layout parameters and
selecting it according to actual soil conditions, and \textcite{Tian2024}
recommend burying the pipe as far above the restrictive/critical layer as
possible when shallow, breaking a shallow restrictive layer by deep ploughing,
and combining methods to enhance drainage. \textcite{Tian2023c} identify
layout-parameter selection as the next research target and suggest amendments
such as organic fertiliser, straw biochar, or low-permeability mulch film, and
\textcite{Zhang2024d} call for adding factors such as buried-pipe depth and
fertiliser ratios to improve the saline-soil improvement system.
\textcite{Zhai2023} advocate optimising hidden-pipe layout through a
multi-scheme decision model and distinguishing nitrous-oxide-emission sources,
and \textcite{Han2023} plan to include construction cost when clarifying the
appropriate range of buried depths. \textcite{Qian2025b} recommend improving
the design drainage objective when using models to set pipe parameters,
\textcite{Ren2021} call for studying combination forms under different rainfall
and unsaturated initial conditions using mathematical models of economic
factors, and \textcite{Zheng2025} propose field experiments on the interaction
between drain spacing and outlet height to develop water- and salt-regulation
control strategies. \textcite{SerranoCoronel2018} state that water-movement
behaviour under different configurations must be studied, and
\textcite{Zegeye2021} suggest modifying drain depth and spacing using better
soil-permeability information and integrating bank dewatering with other gully
rehabilitation measures.

\subsection{Watershed-scale representation, drain mapping, and connectivity}

Scaling drainage processes from plot to catchment raises its own frontier.
\textcite{Muma2016} defer the distributed impact of subsurface drainage on
catchment behaviour to further work, and \textcite{RevueltaAcosta2021} propose
extending the WEPP model to quantify subsurface drainage under controlled water
tables and at watershed scale, and to improve estimates under shallow water
tables during winter. \textcite{Thomas2016} note that describing basin shape
and topographic impacts on tile-drainage contribution to streamflow requires
further research and that calibration and validation for prediction remain to
be investigated. \textcite{Rahman2014} offer their coupling of a watershed
model with streamflow-response analysis as a transferable approach for
evaluating land-use and climate impacts, and \textcite{Schilling2015} identify
the effect of regularly spaced tile drainage on groundwater travel times as an
open question. Streamflow prediction in tile-drained landscapes remains
difficult: \textcite{Wessel2022} state that improving predictions of
subsurface runoff through tile drains is a pressing need, that the validity of
the curve-number approach at different locations must be confirmed, and that
whether a modified curve-number method or a more process-based model is
required must be determined.

Locating and inventorying drainage networks is a prerequisite for many of
these analyses. \textcite{Yimer2023} call for a global study of inter-regional
differences, methods other than interpolation to fill missing values,
approaches to differentiate tile from ditch drains using soil-moisture
attributes, and local groundwater-level measurements in preference to global
datasets. \textcite{Kratt2020} aim to better constrain the optimal
meteorological and soil-water conditions for successful drainage-pipe
detection and to quantify feedbacks between albedo, soil moisture, and heat
transfer to predict optimal data-collection timing, and
\textcite{Allred2018b} note that textural analysis of soil cores is needed to
confirm supposed texture differences, that mapping efficiency could improve
with an all-terrain-vehicle-mounted sled, and that widespread adoption might
yield insight on drainage-practice intensity at small watershed scales.
Surface-ditch connectivity is highlighted as a management-relevant frontier:
\textcite{Moloney2020} argue that mapping surface ditches may be more efficient
than locating individual in-field drains, that future phosphorus-loss and
risk-assessment models should include connectivity metrics via surface ditches,
and that policymakers must consider lag times in nutrient release.
\textcite{Opoku2024} propose applying a bespoke integrated connectivity risk
ranking for phosphorus and nitrogen along open ditches and further assessing
the temporal and spatial variability of soil, weather, drainage system, and
hydrogeochemistry. \textcite{Calsamiglia2018} call for improving connectivity
(IC) algorithms for Mediterranean lowland catchments, continuous monitoring to
establish IC thresholds, and further investigation of thresholds representing
the (dis)activation of connected areas by rainfall frequency and magnitude.

\subsection{Closed depressions, potholes, and edge-of-field practices}

Drained closed depressions and potholes are singled out as poorly understood.
\textcite{Martin2019} argue that the hydrologic behaviour of drained closed
depressions must be better understood to inform conservation practices, that
the frequency, extent, and mechanisms of inundation matter for management, and
that future modelling cannot assume uniformity of potholes at watershed scale.
\textcite{Williams2020} call for research on surface and subsurface hydrologic
processes in drained closed depressions at additional locations, quantification
of the relationship between surface runoff, subsurface tile drainage, and
antecedent soil moisture, and data on the longevity of blind inlets under
no-tillage. \textcite{Feyereisen2015b} report ongoing research on nutrient and
agrochemical fate, alternative nutrient management, the longevity of blind and
gravel inlets, and the effect of a subsurface drainage network on nutrient
losses in closed depressions. Edge-of-field and treatment-system design also
present gaps: \textcite{Gordon2021} suggest raising outlet boards to hold more
water and allow more infiltration, planning for both surface and shallow
groundwater flow, sealing the bottom where groundwater contamination is a
concern, and further studying combined surface/subsurface nitrogen removal with
phosphorus-removing plant harvest; \textcite{Feyereisen2015} call for verifying
design retention times for agricultural residues and incorporating the most
appropriate hydraulic conductivity into bioreactor design models; and
\textcite{Bhattarai2009} note that placement of subsurface drainage beneath
vegetative filter strips needs careful consideration, that uniform surface flow
without preferential paths is essential, and that a filter strip without
subsurface drainage may be effective where rapid water-table lowering is not
required.

\subsection{Climate change, coupled processes, and integrative directions}

Adapting drainage design to a changing climate is an explicit priority.
Climate-change effects and region-specific action plans must be a fundamental
component of drainage design and sizing, and
\textcite{Deelstra2015} tie correct design to understanding weather--soil--
groundwater--discharge interactions under increased precipitation.
\textcite{Tuohy2018} recommend assessing software applicability across a wider
range of Irish drainage systems and soils, quantifying drainage-efficiency
implications for grassland productivity, trafficability, and return on
investment, and using such models to assess climate-change impacts on design
and performance. \textcite{Malota2024} likewise call for assessing model
performance over longer periods and different climates and soils and for
investigating climate-change and soil effects on drainage performance and crop
yields.

Coupled and integrative process representations round out the frontier.
\textcite{Moroizumi2004} note that water-vapour movement and a coupled heat-and-
water-flow model are required for relatively dry soils, and
\textcite{Koivusalo1999} call for a time lag before the steady-state soil-
moisture profile is reached during dry periods and a more consistent
potential-evapotranspiration procedure such as Penman--Monteith.
\textcite{Salo2015} highlight the ability to connect hydraulic functions to
solute variables to enable simulation of density-driven flow, and
\textcite{Varvaris2018} envisage an integrated three-dimensional
hydrobiogeochemical model of solute transport and nitrogen dynamics.
\textcite{Zhang2024} call for addressing intermittent ponding,
evapotranspiration, soil heterogeneity, and three-dimensional flow effects, as
well as uncertainty in the mass-transfer-rate coefficient, and
\textcite{Afruzi2014} note that ponded surface drainage by trapezoidal ditches
warrants further research, while \textcite{Chahar2010} indicate that partially
penetrating ditches require a different analytical solution and that slope
stability may be addressed by refilling ditches with free-draining material.
Model improvement for specific outputs is also flagged: \textcite{Verma2014}
anticipate future model versions enabling more detailed scientific evaluation
of drainage projects, and \textcite{Nousiainen2015} call for further study of
distributed soil hydraulic properties on inefficiently functioning drains.
\textcite{Augeard2005} plan further experiments to confirm winter runoff and
validate suspended-sediment-concentration evolution as a surface-runoff
indicator, supported by hydrological modelling of runoff and pressure heads.

Finally, several authors emphasise that technical advances must be coupled with
socio-economic considerations, extension, and demonstration.
\textcite{BenAissa2013} advocate a participatory approach to design, testing
mixed traditional-ditch and modern-pipe systems and adequate envelopes in situ,
and designing systems for easy farmer inspection and maintenance, echoed by
\textcite{HajAmor2019} on raising farmer and stakeholder awareness of water and
soil management. \textcite{MCFARLANE1992} call for greater extension through
demonstrations and documented case studies and for better research on drains'
effects on salinity, erosion, structural decline, and phosphate export, and
\textcite{Mante2018} note that field managers must weigh the benefit of
narrower spacing against impacts on field operations and long-term yield.
Cost feasibility is a recurring concern---for changing drainage layout, and
\textcite{Imran2021} stresses it specifically for redesigning main drains for
free gravity outflow---and \textcite{Malota2015} suggest linking DRAINMOD
outputs to a drainage-system cost-minimisation model. \textcite{Lamsodis2011}
add operational recommendations, including annual drain-mouth cleaning and
managed control of beaver populations, while \textcite{DarziNaftchali2017}
argue that surface- and groundwater susceptibility to nitrogen losses should be
a major criterion in drain depth and spacing design and that measures are
needed at first planting after installation. \textcite{HARRIS1993} note that
an increased drainage design standard may be necessary to avoid pipe damage and
that direct drilling on sloping undrained land should be avoided, and
\textcite{CAMP1982} stress ensuring adequate surface and internal drainage for
flat-planted sugarcane. Collectively these directions point toward an
integrated research agenda in which improved process representation, denser and
longer-term data, and explicit economic and management context advance
together.